%%%%%%%%%%  scan idx 52  |  folio 39  %%%%%%%%%%
soit bien un anneau gradué résulte du théorème suivant de Chow :

\textbf{Théorème 2.1.}\footnote{Par suite d'un malentendu de la part de l'auteur, le théorème énonce plus qu'il n'est actuellement connu. En fait, on utilise seulement le «~moving lemma~» de Chow, disant qu'on peut bouger deux cycles dans leurs classes de sorte que leur intersection soit non excédentaire~; cf. Séminaire Chevalley 1958, Anneau de Chow et applications.} \emph{Tout cycle sur $X$ est rationnellement équivalent à un cycle non singulier. Étant donnés des éléments de $A^i(X)$ et $A^j(X)$, on peut trouver des représentants non singuliers de ces éléments, soient $Y^{n-i}$ et $Z^{n-j}$, qui se coupent partout transversalement.}

(Un cycle est dit non singulier si ses composantes sont des sous-variétés non singulières. On dit que deux cycles se coupent partout transversalement si toute composante de l'un coupe partout transversalement toute composante de l'autre, i.e. si les deux variétés sécantes sont non singulières aux points d'intersections et si leurs variétés tangentes y sont en position générale \dots)

Soit $f$ un morphisme (i.e. une application régulière) de $X$ dans $Y$ ($X$ et $Y$ étant quasi-projectives et non singulières), alors $f$ définit (en utilisant le théorème précédent) un homomorphisme d'anneaux gradués :
\begin{equation}\tag{2.1}
f^*:A(Y)\longrightarrow A(X)
\end{equation}
(\emph{image réciproque de classes de cycles}). Ainsi $A(X)$ devient un foncteur contravariant en $X$.

Soit $f$ un morphisme \emph{propre} de $X^n$ dans $Y^m$ (i.e. dans la terminologie de Weil, le graphe de $f$ est complet au-dessus de $Y$, cf [3]). Alors $f$ définit aussi un homomorphisme additif
\begin{equation}\tag{2.2}
f_*:A(X^n)\longrightarrow A(Y^m)
\end{equation}
conservant la dimension des cycles, donc augmentant le degré de $m-n$.




% --- Source pages 47--75 ---
